\subsection{映射到乘积的扩展}

定义 2. 设 \((\mathbf{X}_t)_{t \in I}, (\mathbf{Y}_t)_{t \in I}\) 是两个集合族，并且设 \((g_t)_{t \in I}\) 为一个具有相同指标集I的函数族，使得 \(g_t\) 是从 \(\mathbf{X}_t\) 到 \(\mathbf{Y}_t\) 的映射，对于每一个 \(t \in I\) 。对于每一个 \(f \in \prod_{t \in I} \mathbf{X}_t\) ，令 \(u_f\) 为函数 \(t = g_t(f(t))\) ( \(t \in I\) )的图像，它是 \(\prod_{t \in I} \mathbf{Y}_t\) 的一个元素。映射 \(f \to u_f\) （从 \(\prod_{t \in I} \mathbf{X}_t\) 到 \(\prod_{t \in I} \mathbf{Y}_t\) ）被称为映射族 \((g_t)_{t \in I}\) 到乘积的典范扩张（或简称为扩张）；它有时也被称为映射族 \((g_t)_{t \in I}\) 的乘积.

当使用指标记号时，族 \((g_{i})_{i\in I}\) 的乘积是函数 \((x_{i})_{i\in I}\to (g_{i}(x_{i}))_{i\in I}\) ；它有时记作 \((g_{i})_{i\in I}\) .

如果 \(\mathbf{I} = \{\alpha, \beta\}\) ，其中 \(\alpha\) 和 \(\beta\) 是不同的，则映射族 \((g_t)_{t \in \mathbf{I}}\) 到乘积的扩张正是 \(\psi \circ (g_\alpha \times g_\beta) \circ \varphi\) ，其中 \(\varphi\) 表示从 \(\prod_{i \in \mathbf{I}} \mathbf{X}_i\) 到上 \(\mathbf{X}_\alpha \times \mathbf{X}_\beta\) （第3节）的典范映射，且 \(\psi\) 表示从 \(\mathbf{Y}_\alpha \times \mathbf{Y}_\beta\) 到上 \(\prod_{i \in \mathbf{I}} \mathbf{Y}_i\) 的典范映射.

命题11。设 \((\mathbf{X}_t)_{t \in I}, (\mathbf{Y}_t)_{t \in I}, (\mathbf{Z}_t)_{t \in I}\) 是三个集合族，并设 \((g_t)_{t \in I}, (g_t')_{t \in I}\) 是两个函数族，它们具有相同的指标集，使得 \(g_t\) 是从 \(\mathbf{X}_t\) 到 \(\mathbf{Y}_t\) 的映射并且 \(g_t'\) 是从 \(\mathbf{Y}_t\) 到 \(\mathbf{Z}_t\) 的一个映射，对每个 \(t \in I\) . 设 \(g\) 和 \(g'\) 为族 \((g_t)_{t \in I}\) 和 \((g_t')_{t \in I}\) 到乘积的扩张。然后族 \((g_t' \circ g_t)_{t \in I}\) 到乘积的扩张等于 \(g'\circ g\) .

这直接由定义2得出。

推论。设 \((\mathbf{X}_i)_{i\in I}, (\mathbf{Y}_i)_{i\in I}\) 是两个集合族，并设 \((g_i)_{i\in I}\) 是一个函数族。如果 \(g_i\) 是从 \(\mathbf{X}_i\) 到 \(\mathbf{Y}_i\) 的一个单射（相应地，满射），对于每一个 \(i\in I\) ，那么 \(g\) （族 \((g_i)_{i\in I}\) 到乘积的扩张）是从 \(\prod_{i\in I} \mathbf{X}_i\) 到 \(\prod_{i\in I} \mathbf{Y}_i\) 的一个单射（相应地，满射）.

\begin{enumerate}
\def\labelenumi{(\arabic{enumi})}
\item
  让我们假设 \(X_{i} \neq \emptyset\) 对于每一个 \(i \in I\) ；否则结果是平凡的。假设 \(g_{i}\) 对每个 \(i \in I\) ，并且让 \(r_{i}\) 为 \(g_{i}\) （§3，第8号，定义11），因此 \(r_{i} \circ g_{i}\) 是 \(X_{i}\) 。设 r 为族 \((r_{i})_{i \in I}\) ；由于 \(r \circ g\) 是恒等映射族到乘积的扩展 \(I_{X_{i}}\) , \(r \circ g\) 是 \(\prod_{i \in I} X_{i}\) ，因此 g 是单射（§3，命题8）。
\item
  假设 \(g_{i}\) 是……的满射 \(X_{i}\) 到上 \(Y_{i}\) 对于每一个 \(i \in I\) ，并且让 \(s_{i}\) 是……的一个截面 \(g_{i}\) ( \(\S 3\) ，第8号，定义11），因此 \(g_{i} \circ s_{i}\) 是 \(Y_{i}\) 。如果 \(s\) 是族到乘积的扩张 \((s_{i})_{i \in I}\) ，那么 \(g \circ s\) 是恒等映射族到乘积的扩展 \(I_{Y_{i}}\) 并且因此是 \(\prod_{i \in I} Y_{i}\) ；因此 \(g\) 是满射（ \(\S 3\) ，命题8）。
\end{enumerate}

让 \((\mathbf{X}_i)_{i\in I}\) 设为一个集合族，并令 E 为一个集合。对于每个映射 \(f\) 将E映入 \(\prod_{i\in I} \mathbf{X}_i\) , \(\operatorname{pr}_i \circ f\) 是E到 \(\mathbf{X}_i\) 。如果 \(\overline{f}\) 是该映射族到乘积的延拓，且若 \(d\) 是E到 \(\mathbf{E}^{\mathrm{I}}\) 的对角映射，则显然 \(f = \overline{f} \circ d\) 。反过来，设 \((f_{t})_{t \in \mathbf{I}}\) 是这样一个函数族，使得 \(f_{t}\) 是E到 \(\mathbf{X}_{t}\) 对于每一个 \(i \in \mathbf{I}\) ，并且让 \(\overline{f}\) 设为此族到乘积的扩张；于是我们有 \(\mathbf{pr}_{t} \circ (\overline{f} \circ d) = f_{t}\) 对于每一个 \(t \in \mathbf{I}\) . 滥用语言，该映射 \(\overline{f} \circ d\) 也写作 \((f_{t})_{t \in \mathbf{I}}\) 。这样，我们就定义了该集合的一一映射 \(\prod_{i \in \mathbf{I}} \mathbf{X}_{t}^{\mathbf{E}}\) 到集合 \(\left(\prod_{i \in \mathbf{I}} \mathbf{X}_{t}\right)^{\mathbf{E}}\) 此映射及其逆映射被称为典范的。

